% 张三慧第4版第二卷第26章《波粒二象性》原书思考题26.1--26.13
% 主教材V2扫描PDF物理页368；本文件与正式“习题26.1--26.22”分开。
\section*{第26章\quad 波粒二象性：思考题}
\paragraph{思考题 26.1}
霓虹灯的发光与熔炉中炽热铁水的发光分别属于什么机理？是否都可视为单纯热辐射？
\paragraph{思考题 26.2}
人体温度约310K也会向外辐射电磁波，
为什么在黑暗环境中人眼看不见人体自身的热辐射？
\paragraph{思考题 26.3}
新粉刷的房屋即使白墙明亮，
从白天室外远处看，其敞开的窗口仍显得暗甚至黑，
如何用空腔辐射和多重反射解释？
\paragraph{思考题 26.4}
把白瓷砖一半涂上烟煤、在火炉中同时加热，
高温时哪一半更显得明亮？说明光谱发射率与
表面黑度、温度的关系。
\paragraph{思考题 26.5}
洛阳王城公园中的黑牡丹为什么通常需要
避免强烈日照或置于受控环境培育？
从光吸收和热辐射角度提出物理解释，并说明
其不能代替完整的园艺学证据。
\paragraph{思考题 26.6}
假设普朗克常量比现实世界大
\(10^{34}\)倍，如何影响量子能级间距、
热辐射及宏观物体可观测的量子特征？
\paragraph{思考题 26.7}
在光电效应实验中，
(1) 光频率固定、增强入射光强；
(2) 光强固定、提高入射光频率，
分别对光电子的发射数目和最大动能有何影响？
\paragraph{思考题 26.8}
用单色光照射同一金属的洁净表面，
为什么逸出的光电子速率可以不同？
\paragraph{思考题 26.9}
可见光与电子相互作用是否也存在康普顿散射？
为什么通常很难观测到可见光的康普顿波长变化？
\paragraph{思考题 26.10}
为什么分析光电效应常关注单光子的能量守恒，
分析康普顿散射却必须同时使用能量和动量守恒？
\paragraph{思考题 26.11}
电子与质子具有同样的非相对论动能，
哪一种粒子的德布罗意波长更长？
\paragraph{思考题 26.12}
在形式极限 \(h\to0\) 下，
波粒二象性的物质波和经典近似如何改变？
若另作 \(c\to\infty\) 的非相对论极限，
时空相对论性有何变化？
\paragraph{思考题 26.13}
由不确定性原理，一个分子即使处于0K，
其内部运动能否完全停止？
区分局域束缚振动与自由、非局域的平面波态。
